% !TeX root = ../tesis.tex

\chapter{\textit{Finite Integration Technique}}
\label{chapter:FIT}

\vspace*{7em}

%
% --- JAUA ---
% Igual que en pasado: Primero menciona que debes calcular las porpiedades del eritroxito pweo que no se cuenta con una solución analítica, por lo tanto recurres a métpdps niméricos.
%

Debido a la geometría bicóncava de los eritrocitos, su respuesta electromagnética no puede describirse mediante una solución analítica, por lo que se aborda de manera numérica. En este capítulo se presenta la formulación del método de \textit{Finite Integration Technique} (FIT) y su aplicación al estudio del esparcimiento de luz por una esfera. En primer lugar, se describe el método a partir de las ecuaciones de Maxwell en su forma integral y su discretización sobre mallas duales. Se introducen las ecuaciones de malla de Maxwell, junto con las relaciones constitutivas discretas que fundamentan el esquema numérico, así como su formulación en el dominio del tiempo. Posteriormente, se analiza la convergencia de las simulaciones realizadas en \textit{Computer Simulation Technology} (CST) mediante la comparación con la solución analítica de la teoría de Mie \cite{bohrenAbsorptionScatteringLight2008}. Este análisis se lleva a cabo evaluando la influencia de los principales parámetros numéricos —como el tamaño del dominio computacional, el mallado y las condiciones de frontera— sobre la precisión y el costo computacional. A partir de estos resultados, se establecen los criterios de discretización que se emplean como base en las simulaciones presentadas en los capítulos siguientes.




\section{Ecuaciones de malla de Maxwell}
\label{section:MGE}

El método \textit{Finite Integration Technique} (FIT) se basa en la discretización espacial de las ecuaciones de Maxwell a partir de su forma integral, obteniendo así las Ecuaciones de malla de Maxwell \cite{clemensDiscreteElectromagnetismFinite2001}. En este método las magnitudes electromagnéticas se asocian a los objetos geométricos de la malla sobre los cuales se definen dichas integrales \cite{dassaultsystemesCSTMicrowaveStudio2003}.

El primer paso del proceso de discretización empleado en FIT consiste en definir el mallado de integración. Para ello, se emplea una discretización tipo Yee \cite{yeeNumericalSolutionInitial1966}. En este contexto, el problema electromagnético ---generalmente formulado en una región con frontera abierta---  se restringe a un dominio computacional finito $\Omega\in\mathbb{R}^3$. Posteriormente, dicho dominio se descompone en un número finito de subregiones $V_i$, tales como tetraedros o hexaedros, que encajan exactamente entre sí. Esto implica que la intersección entre dos celdas distintas  es vacía o corresponde a un elemento geométrico compartido de menor dimensión (una cara, una arista o un vértice). Esta descomposición da lugar a un sistema coordenado ortogonal $G$, que constituye la malla computacional del problema~\cite{clemensDiscreteElectromagnetismFinite2001}. Para simplificar el desarrollo del método, se asume que $\Omega$ tiene forma paralelepipédica y que su discretización se realiza mediante una malla cartesiana obtenida como producto tensorial de subdivisiones unidimensionales, de modo que \cite{hadabSolutionMaxwellsEquation2023}
%
\begin{align}
	G=\{(x_i,y_j,z_k); i\in\{1,...,I\};j\in\{1,...,J\};k\in\{1,...,K\}\}.
\end{align}
%
Los nodos de la malla $(x_i,y_j,z_k)$ se indexan mediante los índices $i,j,k$ a lo largo de los ejes $x,y,z$. De esta forma, el mallado contiene un total de $N_p=IJK$ nodos y $(I-1)(J-1)(K-1)$ celdas \cite{hadabSolutionMaxwellsEquation2023}. Para facilitar la implementación numérica, los nodos se enumeran mediante un único índice $n$ definido como \cite{hadabSolutionMaxwellsEquation2023}
%
\begin{equation}
	n=1+(i-1)M_x+(j-1)M_y+(k-1)M_z,
	\label{eq:index_system}
\end{equation}
% 
donde $M_x=1,M_y=I$ y $M_z=IJ$. En la Fig. \ref{fig:G_complex} se muestra el sistema coordenado $G$ para un conjunto de celdas cartesianas. Asimismo, se muestran las componentes de los nodos en los extremos de la malla, así como la identificación del volumen $V_n$ y el área $S_{u,m}$ asociados a los elementos vinculados a los nodos $n$ y $m$, respectivamente.
%
\begin{figure}[h!]
	\centering
	\includegraphics[width=8cm]{3-FIT/G_complex.pdf}
	\caption{Representación de una malla tridimensional de tipo Yee ($G$) en un sistema de coordenadas cartesiano. Se muestra la discretización del volumen de interés donde cada vértice actúa como un nodo $n$, con $N = I \times J \times K$ nodos totales vinculados mediante el sistema de indexación dado por la Ec. \eqref{eq:index_system}. }
	\label{fig:G_complex}
\end{figure}
%


Una vez definido $G$, el método FIT se formula a partir de las ecuaciones de Maxwell en su forma integral que, en unidades del SI, están dadas por \cite{griffithsIntroductionElectrodynamics2023}

%
\begin{subequations} \label{eqs:Maxwell_int}
	\begin{tcolorbox}[
		ams align, breakable]
		\oint_{S}\vb{D} \cdot \dd{\vb{a}}&=Q_{\text{ext}}, &\mbox{(Ley de Gauss eléctrica)}  
		\label{seq:GE_int} \\
		\oint_{S}\vb{B} \cdot \dd{\vb{a}} &= 0,						&\mbox{(Ley de Gauss magnética)}   
		\label{seq:GM_int} \\
		\oint_{\partial S}\vb{E} \cdot d\vb{l} &= -\frac{\text{d}}{\dd{t}}\int_{S}\vb{B}\cdot\dd{\vb{a}}, 	&\mbox{(Ley de Faraday-Lenz)}		
		\label{seq:FL_int}\\
		\oint_{\partial S}\vb{H} \cdot d\vb{l} &= \int_{S}\left(\vb{J}_{\text{ext}}+\frac{\partial \vb{D}}{\partial t}\right)\cdot\dd{\vb{a}}, &
		\mbox{(Ley de Ampère-Maxwell)} \label{seq:AM_int}
\end{tcolorbox}\end{subequations}\noindent
%
donde se obvian las dependencias espaciales y temporales; $\rho_{\text{ext}}$ es la densidad volumétrica de carga externa y  $\vb{J}_{\text{ext}}$ es la densidad volumétrica de corriente  externa. A partir de las ecuaciones de Maxwell en su forma integral, se construye la discretización del método FIT. En particular, al aplicar la ley de Faraday–Lenz [Ec. \ref{seq:FL_int}] sobre una cara $S_{z,n}$ asociada a una celda con volumen $V_n$ se obtiene~\cite{hadabSolutionMaxwellsEquation2023}
%
\begin{equation}
	\oint_{\partial S}\vb{E} \cdot d\vb{l}=(x_{i+1}-x_i)E_{x,n}+(y_{j+1}-y_j)E_{y,n+1}+(x_{i+1}-x_i)E_{u,n+I}+(y_{j+1}-y_j)E_{y,n},
	\label{eq:lineint_Yee}
\end{equation}
%
donde la integral de línea se evalúa sobre los cuatro segmentos que conforman el contorno de la cara $S_{z,n}$. Por consiguiente, la integral de área considera el producto de los segmentos 
%
\begin{equation}
	\int_{S}\vb{B}\cdot\dd{\vb{a}}=B_{z,n}(x_{i+1}-x_i)(y_{j+1}-y_j),
	\label{eq:surfint_Yee}
\end{equation}
%
de modo que la Ec. \eqref{seq:FL_int} se reescribe como %
%
\begin{align}
	(x_{i+1}-x_i)E_{x,n}+(y_{j+1}-y_j)E_{y,n+1}&+(x_{i+1}-x_i)E_{u,n+I}+ \nonumber\\
	&+(y_{j+1}-y_j)E_{y,n}
	=-\frac{\dd}{\dd{t}}B_{z,n}(x_{i+1}-x_i)(y_{j+1}-y_j).
	\label{eq:Faraday_grid}
\end{align}
%
En la Fig.~\ref{subfig:FIT_cell_electric} se muestra la superficie $S_{z,n}$ de una celda en la cual los vectores en rojo y azul oscuro representan, respectivamente, la circulación del campo eléctrico a lo largo de las aristas y el flujo magnético que atraviesa la superficie. 
%
\begin{figure}[tb]
	\centering
	\sidesubfloat[First image]{\hspace{-0.8cm}{
			\includegraphics[width=7.8cm]{3-FIT/FIT_cell.pdf}\label{subfig:FIT_cell_electric}}}\hspace{0.1cm}
	\sidesubfloat[Second image]{\hspace{-0.5cm}{	\includegraphics[width=7.5cm]{3-FIT/FIT_cell_fluxB.pdf}\label{subfig:FIT_cell_magnetic}}}\vspace{-0.1cm}
	\caption{\textbf{a)} Cara $S_{z,n}$ de la celda $V_{n}$ del complejo $G$ en la que las flechas rojas indican la circulación del campo eléctrico a lo largo de las aristas que conforman el contorno de la cara, mientras que la flecha azul representa el flujo del campo magnético a través de dicha superficie, asociado a $B_{z,n}$. \textbf{b)} Celda $V_{n}$ del complejo $G$ que muestra el flujo de campo magnético que atraviesa cada una de sus superficies. Los subíndices asociados a cada flujo corresponden a la superficie de la cara que atraviesan.}
	\label{fig:FIT_cell}
\end{figure}
%

La formulación anterior es válida para cualquier cara de $G$ y puede extenderse a superficies formadas por la unión de varias caras, ya que la suma de las circulaciones sobre los contornos individuales equivale a la circulación sobre el contorno total de la superficie \cite{arfkenMathematicalMethodsPhysicists2011}. Con el fin de extender esta relación a todo el dominio, se introduce una formulación matricial que permite representar de manera compacta tanto las magnitudes electromagnéticas como los operadores integrales discretos.
En este marco, los campos electromagnéticos se organizan como vectores de dimensión $3N$. En particular, los campos $\vb{E}, \vb{D}, \vb{B}$ y $\vb{H}$ se representan mediante los vectores $\vb{e}, \vb{d}, \vb{b}$ y $\vb{h}$, respectivamente, definidos como \cite{hadabSolutionMaxwellsEquation2023}
\begin{equation}
	\vb{e}=\begin{bmatrix}
		E_{x,1}\\
		E_{x,2}\\
		\vdots\\
		E_{z,N} 
	\end{bmatrix},\;\;
	\vb{d}=\begin{bmatrix}
		D_{x,1}\\
		D_{x,2}\\
		\vdots\\
		D_{z,N} 
	\end{bmatrix},\;\;
	\vb{b}=\begin{bmatrix}
		B_{x,1}\\
		B_{x,2}\\
		\vdots\\
		B_{z,N} 
	\end{bmatrix},\;\;
	\vb{h}=\begin{bmatrix}
		H_{x,1}\\
		H_{x,2}\\
		\vdots\\
		H_{z,N} 
	\end{bmatrix},
\end{equation}
%
y de manera análoga se definen los vectores $\vb{j}$ y $\vb{q}$, correspondientes a la discretización de $\vb{J}_{\text{ext}}$ y $Q_{\text{ext}}$, respectivamente. Para construir los operadores integrales discretos, se introducen las matrices $P_u$, con $u=\{x,y,z\}$, que codifican la conectividad de la malla. Estas matrices se definen de forma general como \cite{hadabSolutionMaxwellsEquation2023}
%
\begin{equation}
	(P_u)_{m,p}=
	\begin{cases}
		-1, & \text{si } p=m,\\
		1, & \text{si } p=m+M_u,\\
		0, & \text{en cualquier otro caso},
	\end{cases}
	\qquad (m,p)\in \{1,\dots,N\}^2.
\end{equation}
A partir de estas, se construye el operador rotacional $\mathbb{C}\in\mathbb{R}^{3N\times 3N}$, dado por \cite{hadabSolutionMaxwellsEquation2023}
%
\begin{equation}
	\mathbb{C}=\begin{pmatrix}
		0 & -P_z & P_y\\
		P_z & 0 & -P_x\\
		-P_y & P_x & 0
	\end{pmatrix}.
\end{equation}
Adicionalmente, se introducen matrices diagonales que contienen la información geométrica del mallado. En particular, la matriz $\mathbb{D}_l$ agrupa las longitudes asociadas a las aristas, mientras que la matriz $\mathbb{D}_S$ contiene las áreas de las caras. La primera se define como \cite{hadabSolutionMaxwellsEquation2023}
%
\begin{equation}
	\mathbb{D}_l=\text{diag}(\Delta x_{1},\dotso,\Delta x_{N},\Delta y_{1},\dotso,\Delta y_{N},\Delta z_{1},\dotso,\Delta z_{N}),
\end{equation}
%
donde las longitudes asociadas a cada dirección se definen de manera general como \cite{hadabSolutionMaxwellsEquation2023}
%
\begin{equation}
	\Delta u_n =
	\begin{cases}
		\displaystyle \int_{u_\ell}^{u_{\ell+1}} \left\lVert \frac{\partial \vb{r}}{\partial u} \right\rVert \dd{u}, & \text{si } \ell < N_u,\\
		0, & \text{si } \ell = N_u,
	\end{cases}
\end{equation}
%
con $\ell=\{i,j,k\}$ tal que $u_\ell$ es la coordenada correspondiente ($x_i$, $y_j$ o $z_k$), y $N_u$ el número total de nodos en la dirección $u$ (es decir, $N_x=I$, $N_y=J$, $N_z=K$) \cite{hadabSolutionMaxwellsEquation2023}. De manera análoga, la matriz de áreas está dada por \cite{hadabSolutionMaxwellsEquation2023}
%
\begin{equation}
	\mathbb{D}_S=\text{diag}(S_{x,1},\dotso,S_{x,N}, S_{y,1},\dotso,S_{y,N},S_{z,1},\dotso,S_{z,N}),
\end{equation}
%
con $S_{x,n}=\Delta y_n \Delta z_n, S_{y,n}=\Delta x_n \Delta z_n$ y $S_{z,n}=\Delta x_n \Delta y_n$ \cite{hadabSolutionMaxwellsEquation2023}. Finalmente, la ley de Faraday–Lenz (Ec. \eqref{seq:FL}) empleando la discretización tipo Yee se reescribe en forma matricial como \cite{hadabSolutionMaxwellsEquation2023}
%
\begin{equation}
	\mathbb{C}\mathbb{D}_l\vb{e}=-\mathbb{D}_S\frac{\partial}{\partial t}\vb{b}.
\end{equation}
%

El segundo operador a discretizar es el operador de la divergencia. Para ello, se parte de la Ec. \eqref{seq:GM_int} aplicada a una celda de volumen $V_n$, como se muestra en la Fig.~\ref{fig:FIT_cell}, donde las flechas azules representan los flujos magnéticos asociados a cada una de las caras de la celda. Al aplicar la ley de Gauss magnética siguiendo el esquema de Yee (Ec. \ref{eq:surfint_Yee}), la ecuación integral puede reescribirse como una suma de los flujos magnéticos que atraviesan las seis caras de la celda, dando lugar a~\cite{hadabSolutionMaxwellsEquation2023}
%
\begin{align}
	B_{z,n}(x_{i+1}-x_i)(y_{j+1}-y_j)
	&-B_{z,n+1}(x_{i+1}-x_i)(y_{j+1}-y_j)\nonumber\\
	&+B_{x,n}(y_{j+1}-y_j)(z_{k+1}-z_k)\nonumber\\
	&-B_{x,n+1}(y_{j+1}-y_j)(z_{k+1}-z_k)\nonumber\\
	&+B_{y,n}(x_{i+1}-x_i)(z_{k+1}-z_k)\nonumber\\
	&-B_{y,n+1}(x_{i+1}-x_i)(z_{k+1}-z_k)=0.
	\label{eq:GM-Yee}
\end{align}
%
La Ec. \eqref{eq:GM-Yee} puede escribirse en forma matricial para extenderse en todo el dominio mediante la introducción del operador de divergencia \cite{hadabSolutionMaxwellsEquation2023}
%
\begin{equation}
	\mathbb{S} = \begin{pmatrix} P_x & P_y & P_z \end{pmatrix},
\end{equation}
%
el cual actúa sobre el vector de campo magnético ponderado por las áreas de las caras. De este modo, la ecuación discreta de la ley de Gauss magnética se expresa como \cite{hadabSolutionMaxwellsEquation2023}
%
\begin{equation}
	\mathbb{S}\mathbb{D}_S\vb{b} = 0.
\end{equation}
%%

La discretización de las ecuaciones de Maxwell restantes requiere la introducción de un segundo sistema coordenado $\tilde{G}$, dual al complejo primario $G$ \cite{hadabSolutionMaxwellsEquation2023}
%
\begin{align}
	\tilde{G}=\{(\tilde{x}_i,\tilde{y}_j,\tilde{z}_k); i\in\{1,...,I\};j\in\{1,...,J\};k\in\{1,...,K\}\}.
\end{align}
%
 Este complejo dual se construye tomando los centros de las celdas de $G$ como vértices de las celdas de $\tilde{G}$. Con esta definición se obtiene una correspondencia uno a uno entre las aristas de $G$ y las caras de $\tilde{G}$, y viceversa, con la malla dual desplazada media celda respecto a la malla primaria \cite{clemensDiscreteElectromagnetismFinite2001}. En consecuencia, los nodos de la malla dual ($\tilde{x}_i,\tilde{y}_j,\tilde{z}_k$) se definen como los puntos medios entre nodos consecutivos de la malla primaria \cite{hadabSolutionMaxwellsEquation2023}
 %
 \begin{equation}
 	 \tilde{u}_{\ell} =\dfrac{u_{l+1}+u_l}{2}, \begin{cases}
 		 l \in \{1,...,N_u-1\}\\
 		\tilde{u}_{N_u}=0
 	\end{cases}.
 \end{equation}
 %
 En la Fig. \ref{fig:FIT_cell_dual} se muestra la relación de ortogonalidad entre una celda del complejo primario $G$ (celda amarilla) y su correspondiente celda dual $\tilde{G}$ (celda gris). Además, se muestran las cantidades asociadas a cada complejo ($\vb{E}$ y $\vb{B}$ asociados a $G$, mientras que $\vb{D}$ y $\vb{H}$ asociados a $\tilde{G}$) y su localización espacial en cada celda.
 
 %
 \begin{figure}[h!]
 	\centering
 	\includegraphics[width=8cm]{3-FIT/FIT_dual_grid.pdf}
 	\caption{Celda $V_{n}$ del complejo $G$ y celda dual $\tilde{V}_{n}$ del complejo $\tilde{G}$ desplazada media celda en todas las direcciones y ortogonal a la primera. Se muestran las cantidades asociadas a una cara $S_{z,n}$ y $\tilde{S}_{\tilde{z},n}$ del complejo $G$ y $\tilde{G}$, respectivamente. Se muestran en azul las cantidades relacionadas con los campos magnéticos y en rojo con los campos eléctricos.}
 	\label{fig:FIT_cell_dual}
 \end{figure}
 %
 \noindent De manera análoga al caso del complejo primario, se introducen matrices diagonales que contienen la información geométrica de la malla dual. En particular, la matriz dual de longitudes de integración  $\tilde{\mathbb{D}}_{\ell}$ se define como \cite{hadabSolutionMaxwellsEquation2023}
 %
 \begin{equation}
 	 \tilde{\mathbb{D}}_{\ell}=\text{diag}(\Delta\tilde{x}_{1},\dotso,\Delta\tilde{x}_{N},\Delta\tilde{y}_{1},\dotso,\Delta\tilde{y}_{N},\Delta\tilde{z}_{1},\dotso,\Delta\tilde{z}_{N}),
 \end{equation}
 %
 donde las longitudes asociadas a cada dirección se definen de manera general como \cite{hadabSolutionMaxwellsEquation2023}
 %
 \begin{equation}
 	\Delta \tilde{u}_n =
 	\begin{cases}
 		\dfrac{\Delta u_n+ \Delta u_{n-M_u}}{2}, &  2<\ell < N_u-1,\\
 		\dfrac{\Delta u_n}{2}, & \ell = 1,\\
 		\dfrac{\Delta u_{n-M_u}}{2}, & \ell=M_u.
 	\end{cases}
 \end{equation}
 %
De forma similar, la matriz de áreas de la malla dual $\tilde{\mathbb{D}}_S$ está dada por \cite{hadabSolutionMaxwellsEquation2023}
 %
 \begin{equation}
 	\mathbb{D}_S=\text{diag}(\tilde{S}_{x,1},\dotso,\tilde{S}_{x,N}, \tilde{S}_{y,1},\dotso,\tilde{S}_{y,N},\tilde{S}_{z,1},\dotso,\tilde{S}_{z,N}),
 \end{equation}
 %
 con $\tilde{S}_{x,n}=\Delta \tilde{y}_n \Delta \tilde{z}_n, \tilde{S}_{y,n}=\Delta \tilde{x}_n \Delta \tilde{z}_n$ y $\tilde{S}_{z,n}=\Delta \tilde{x}_n \Delta \tilde{y}_n$ \cite{hadabSolutionMaxwellsEquation2023}. 
 
 Los operadores integrales discretos en la malla dual se relacionan directamente con los del complejo primario. En particular, el operador rotacional dual $\tilde{\mathbb{C}}$ y el operador divergencia dual $\tilde{\mathbb{S}}$ están dados por \cite{hadabSolutionMaxwellsEquation2023}
 %
 \begin{align}
 	\tilde{\mathbb{C}}&=\mathbb{C}^T\\
 	\tilde{\mathbb{S}}&=(P_x^T,P_y^T,P_z^T).
 \end{align}
 %
 De esta forma, para el par de sistemas $\{G,\tilde{G}\}$, el conjunto completo de ecuaciones matriciales discretas —conocidas como \textit{Maxwell Grid Equations} (MGE)— se escribe como \cite{hadabSolutionMaxwellsEquation2023}
%
\begin{subequations} \label{eqs:MGE}
	\begin{tcolorbox}[
		ams align, breakable]
	\tilde{\mathbb{S}}\tilde{\mathbb{D}}_S \vb{d}&= \vb{q}, &\mbox{(Ley de Gauss eléctrica)}  
		\label{seq:GE_MGE} \\
	\mathbb{S}\mathbb{D}_S\vb{b}&=0,						&\mbox{(Ley de Gauss magnética)}   
		\label{seq:GM_MGE} \\
		\mathbb{C}\mathbb{D}_l\vb{e}&=-\mathbb{D}_S\frac{\partial}{\partial t}\vb{b},	&\mbox{(Ley de Faraday-Lenz)}		
		\label{seq:FL_MGE}\\
		\tilde{\mathbb{C}}\tilde{\mathbb{D}}_l\vb{h} &= \tilde{\mathbb{D}}_S \left( \frac{\text{d}}{\dd{t}}\vb{d}+\vb{j}\right). &
		\mbox{(Ley de Ampère-Maxwell)} \label{seq:AM_MGE}
\end{tcolorbox}\end{subequations}\noindent
%

Además de las MGE, el método FIT incorpora las relaciones constitutivas discretas, que describen la respuesta del material. Estas relaciones se escriben en forma matricial como \cite{hadabSolutionMaxwellsEquation2023}
%
\begin{equation}
	\vb{d}=\tilde{\mathbb{D}}_{\varepsilon}\vb{e},\;\;\; \text{y}\;\;\;\vb{b}=\mathbb{D}_{\mu}\vb{h}
\end{equation}
%
donde $\tilde{\mathbb{D}}_{\varepsilon}$  y $\mathbb{D}_{\mu}$ son las matrices discretas de permitividad y la permeabilidad, respectivamente. Para medios isótropos, estas matrices son diagonales \cite{clemensDiscreteElectromagnetismFinite2001}.

